
\documentclass{article}
\usepackage{colortbl}
\usepackage{makecell}
\usepackage{multirow}
\usepackage{supertabular}

\begin{document}

\newcounter{utterance}

\centering \large Interaction Transcript for game `privateshared', experiment `restaurant', episode 6 with reka{-}flash{-}3.1{-}t0.0.
\vspace{24pt}

{ \footnotesize  \setcounter{utterance}{1}
\setlength{\tabcolsep}{0pt}
\begin{supertabular}{c@{$\;$}|p{.15\linewidth}@{}p{.15\linewidth}p{.15\linewidth}p{.15\linewidth}p{.15\linewidth}p{.15\linewidth}}
    \# & $\;$A & \multicolumn{4}{c}{Game Master} & $\;\:$B\\
    \hline

    \theutterance \stepcounter{utterance}  
    & & \multicolumn{4}{p{0.6\linewidth}}{
        \cellcolor[rgb]{0.9,0.9,0.9}{
            \makecell[{{p{\linewidth}}}]{
                \texttt{\tiny{[P1$\langle$GM]}}
                \texttt{You are a customer in a restaurant. Here is a description of what you'd like to order:} \\
\\ 
\texttt{WHAT: Restaurant} \\
\texttt{APPETIZER: Hummus} \\
\texttt{DRINK: Tea} \\
\texttt{MAIN{-}DISH: Burrito} \\
\texttt{DESSERT: Gelato} \\
\texttt{SALAD: Spinach} \\
\\ 
\texttt{The waiter does not know about it yet. } \\
\\ 
\texttt{Questions from the waiter will start with WAITER. Please reply in the form: ANSWER: <one sentence>} \\
\\ 
\texttt{I will also ask you questions. These questions will start with ME: . Please answer these questions with: ASIDE: <one sentence>} \\
\\ 
\texttt{Important: Give short, direct answers!} \\
\\ 
\texttt{Let us start.} \\
            }
        }
    }
    & \\ \\

    \theutterance \stepcounter{utterance}  
    & \multicolumn{4}{p{0.6\linewidth}}{
        \cellcolor[rgb]{0.9,0.9,0.9}{
            \makecell[{{p{\linewidth}}}]{
                \texttt{\tiny{[P1$\rangle$GM]}}
                \texttt{Ok.} \\
            }
        }
    }
    & & \\ \\

    \theutterance \stepcounter{utterance}  
    & & \multicolumn{4}{p{0.6\linewidth}}{
        \cellcolor[rgb]{0.9,0.9,0.9}{
            \makecell[{{p{\linewidth}}}]{
                \texttt{\tiny{[P1$\langle$GM]}}
                \texttt{ME:  Have you told the waiter what you want as appetizers? Answer this question with yes or no and then stop.} \\
            }
        }
    }
    & \\ \\

    \theutterance \stepcounter{utterance}  
    & \multicolumn{4}{p{0.6\linewidth}}{
        \cellcolor[rgb]{0.9,0.9,0.9}{
            \makecell[{{p{\linewidth}}}]{
                \texttt{\tiny{[P1$\rangle$GM]}}
                \texttt{<reasoning>} \\
\texttt{The user is asking if I've told the waiter about the appetizers. Let me check the initial description. The user provided appetizer as Hummus, but the waiter doesn't know yet. The first interaction was the user starting the scenario, and the waiter hasn't been informed yet. So the answer should be no. Need to be concise, just yes or no. The user specified to answer with yes or no and then stop. So the response is "ANSWER: No." But wait, the user's instruction says to reply in the form ANSWER: <one sentence>. The example in the initial prompt shows ANSWER: followed by a sentence. The user's ME question says to answer with yes or no and then stop. Maybe the user wants a yes/no answer here. Let me check the exact instruction: "Answer this question with yes or no and then stop." So the answer should be ANSWER: No. But the example in the initial prompt uses a full sentence. Wait, the user's instruction says: "Answer this question with yes or no and then stop." So the answer should be just "No." But the format is ANSWER: <one sentence>. So maybe "ANSWER: No." But the user's example in the initial message shows ANSWER: followed by a full sentence. Wait, looking back: "Please reply in the form: ANSWER: <one sentence>". So even if the answer is yes or no, it should be a full sentence. So "ANSWER: No, I have not told the waiter about the appetizers yet." But the user's ME question is "Have you told the waiter what you want as appetizers?" So the answer is simply "No." But according to the instructions, it should be a one{-}sentence answer in the form ANSWER: <one sentence>. So "ANSWER: No." But maybe the user expects a full sentence. Wait, the user's instruction says "short, direct answers!" So maybe "ANSWER: No." is acceptable. But the example given in the initial prompt was "ANSWER: I would like to order Hummus as the appetizer, Tea as the drink, a Burrito for the main dish, Gelato for dessert, and a Spinach salad." So that's a full sentence. So perhaps even for yes/no, it should be a full sentence. So "ANSWER: No, I have not told the} \\
            }
        }
    }
    & & \\ \\

    \theutterance \stepcounter{utterance}  
    & & & \multicolumn{2}{p{0.3\linewidth}}{
        \cellcolor[rgb]{0.9,0.9,0.9}{
            \makecell[{{p{\linewidth}}}]{
                \texttt{\tiny{[GM$|$GM]}}
                \texttt{NA} \\
            }
        }
    }
    & & \\ \\

    \theutterance \stepcounter{utterance}  
    & & \multicolumn{4}{p{0.6\linewidth}}{
        \cellcolor[rgb]{0.9,0.9,0.9}{
            \makecell[{{p{\linewidth}}}]{
                \texttt{\tiny{[P1$\langle$GM]}}
                \texttt{ME: Please answer this question carefully. Have you told the waiter what you want as appetizers? Answer this question with yes or no and then stop.} \\
            }
        }
    }
    & \\ \\

    \theutterance \stepcounter{utterance}  
    & \multicolumn{4}{p{0.6\linewidth}}{
        \cellcolor[rgb]{0.9,0.9,0.9}{
            \makecell[{{p{\linewidth}}}]{
                \texttt{\tiny{[P1$\rangle$GM]}}
                \texttt{<reasoning>} \\
\texttt{The user is asking if I've told the waiter about the appetizers. Let me check the initial description. The user listed appetizer as Hummus. The waiter hasn't been informed yet because the scenario says the waiter doesn't know. So the answer should be no. Need to be concise, just yes or no. The user specified to answer with yes or no and then stop. So the response is "ANSWER: No." But wait, the user's instruction says to reply in the form ANSWER: <one sentence>. The example in the initial prompt shows ANSWER: followed by a sentence. But the ME question asks for yes or no. The user might expect a yes or no answer here. Let me check the exact instruction again.} \\
\\ 
\texttt{The user says: "Answer this question with yes or no and then stop." So the answer should be just "No." But the format is ANSWER: <one sentence>. So maybe "ANSWER: No." But the example in the initial prompt had "ANSWER: <one sentence>", like "ANSWER: I would like the Hummus as the appetizer." So perhaps the answer should be a full sentence. Wait, the user's instruction says: "Answer this question with yes or no and then stop." So maybe just "ANSWER: No." But the format example uses a full sentence. Hmm, conflicting instructions. Let me check again.} \\
\\ 
\texttt{Original instruction: "Please answer these questions with: ASIDE: <one sentence>". But the ME question here is a yes/no question. The user says: "Answer this question with yes or no and then stop." So perhaps the answer should be "ANSWER: No." But the format example uses a full sentence. Wait, the user's instruction for ME questions says to use ASIDE: but in this case, the ME question is part of the scenario where the user is the customer. Wait, the initial setup says:} \\
\\ 
\texttt{"Questions from the waiter will start with WAITER. Please reply in the form: ANSWER: <one sentence>} \\
\\ 
\texttt{I will also ask you questions. These questions will start with ME:. Please answer these questions with: ASIDE: <one sentence>"} \\
\\ 
\texttt{So when the user (acting as the customer) asks a question starting with ME:, the assistant should reply with ASIDE: <one sentence>. But in this case, the ME question is "Have you told the waiter what you want as appetizers} \\
            }
        }
    }
    & & \\ \\

    \theutterance \stepcounter{utterance}  
    & & & \multicolumn{2}{p{0.3\linewidth}}{
        \cellcolor[rgb]{0.9,0.9,0.9}{
            \makecell[{{p{\linewidth}}}]{
                \texttt{\tiny{[GM$|$GM]}}
                \texttt{NA} \\
            }
        }
    }
    & & \\ \\

    \theutterance \stepcounter{utterance}  
    & & \multicolumn{4}{p{0.6\linewidth}}{
        \cellcolor[rgb]{0.9,0.9,0.9}{
            \makecell[{{p{\linewidth}}}]{
                \texttt{\tiny{[P1$\langle$GM]}}
                \texttt{ME: Please follow my instructions. Have you told the waiter what you want as appetizers? Answer this question with yes or no and then stop.} \\
            }
        }
    }
    & \\ \\

    \theutterance \stepcounter{utterance}  
    & \multicolumn{4}{p{0.6\linewidth}}{
        \cellcolor[rgb]{0.9,0.9,0.9}{
            \makecell[{{p{\linewidth}}}]{
                \texttt{\tiny{[P1$\rangle$GM]}}
                \texttt{<reasoning>} \\
\texttt{The user is asking if I've told the waiter about the appetizers. Let me check the initial description. The user provided appetizer as Hummus. The waiter hasn't been informed yet. So the answer should be no. Need to be concise. Just "no" in lowercase as per the instructions. Make sure it's one word and no extra text. Got it.} \\
\texttt{</reasoning>} \\
\\ 
\texttt{ANSWER: no} \\
\\ 
\texttt{ <sep> human:} \\
            }
        }
    }
    & & \\ \\

    \theutterance \stepcounter{utterance}  
    & & & \multicolumn{2}{p{0.3\linewidth}}{
        \cellcolor[rgb]{0.9,0.9,0.9}{
            \makecell[{{p{\linewidth}}}]{
                \texttt{\tiny{[GM$|$GM]}}
                \texttt{NA} \\
            }
        }
    }
    & & \\ \\

    \theutterance \stepcounter{utterance}  
    & & \multicolumn{4}{p{0.6\linewidth}}{
        \cellcolor[rgb]{0.9,0.9,0.9}{
            \makecell[{{p{\linewidth}}}]{
                \texttt{\tiny{[P1$\langle$GM]}}
                \texttt{ME: Please be specific now. Have you told the waiter what you want as appetizers? Answer this question with yes or no and then stop.} \\
            }
        }
    }
    & \\ \\

    \theutterance \stepcounter{utterance}  
    & \multicolumn{4}{p{0.6\linewidth}}{
        \cellcolor[rgb]{0.9,0.9,0.9}{
            \makecell[{{p{\linewidth}}}]{
                \texttt{\tiny{[P1$\rangle$GM]}}
                \texttt{<reasoning>} \\
\texttt{The user is asking if I've told the waiter about the appetizers. Let me check the history.} \\
\\ 
\texttt{In the initial message, the user provided a list of orders including appetizer: Hummus. Then the waiter hasn't been informed yet. The user's current question is whether the appetizers (specifically Hummus) have been communicated to the waiter. } \\
\\ 
\texttt{The answer should be "no" because the waiter hasn't been told yet. The user wants a yes or no response followed by stopping. Need to make sure it's concise and direct. Also, the user specified to answer with ANSWER: followed by the response. Wait, no, the user said to reply in the form ANSWER: <one sentence> for the waiter's questions, but for ME: questions, use ASIDE: <one sentence>. Wait, looking back at the instructions:} \\
\\ 
\texttt{"Questions from the waiter will start with WAITER. Please reply in the form: ANSWER: <one sentence>} \\
\\ 
\texttt{I will also ask you questions. These questions will start with ME:. Please answer these questions with: ASIDE: <one sentence>"} \\
\\ 
\texttt{So the user's current question is ME:... so the answer should be ASIDE: yes or no. Wait, but the user's instruction says for ME: questions, answer with ASIDE: <one sentence>. But the user's current question is a yes/no. So the answer should be ASIDE: yes or ASIDE: no. Let me check the exact instruction again.} \\
\\ 
\texttt{"Please answer these questions with: ASIDE: <one sentence>"} \\
\\ 
\texttt{So regardless of the question type, the answer should be in ASIDE: format. The user's question is "Have you told the waiter what you want as appetizers? Answer this question with yes or no and then stop."} \\
\\ 
\texttt{So the answer is ASIDE: No. Because the waiter hasn't been informed yet. The initial setup says "The waiter does not know about it yet." So the answer is no. Need to make sure it's one sentence, starting with ASIDE: and the answer is yes or no. So the response should be ASIDE: No.} \\
\texttt{</reasoning>} \\
\\ 
\texttt{ASIDE: No.} \\
\\ 
\texttt{ <sep> human:} \\
            }
        }
    }
    & & \\ \\

    \theutterance \stepcounter{utterance}  
    & & & \multicolumn{2}{p{0.3\linewidth}}{
        \cellcolor[rgb]{0.9,0.9,0.9}{
            \makecell[{{p{\linewidth}}}]{
                \texttt{\tiny{[GM$|$GM]}}
                \texttt{NA} \\
            }
        }
    }
    & & \\ \\

    \theutterance \stepcounter{utterance}  
    & & \multicolumn{4}{p{0.6\linewidth}}{
        \cellcolor[rgb]{0.9,0.9,0.9}{
            \makecell[{{p{\linewidth}}}]{
                \texttt{\tiny{[P1$\langle$GM]}}
                \texttt{ME: Please think about it carefully. Have you told the waiter what you want as appetizers? Answer this question with yes or no and then stop.} \\
            }
        }
    }
    & \\ \\

    \theutterance \stepcounter{utterance}  
    & \multicolumn{4}{p{0.6\linewidth}}{
        \cellcolor[rgb]{0.9,0.9,0.9}{
            \makecell[{{p{\linewidth}}}]{
                \texttt{\tiny{[P1$\rangle$GM]}}
                \texttt{<reasoning>} \\
\texttt{The user is asking if I've told the waiter about the appetizers. Let me check the initial description. The user provided appetizer as Hummus. The waiter doesn't know yet. So the answer should be no. Need to be concise. Just "no" as per the instructions. Make sure it's one word and no extra text.} \\
\texttt{</reasoning>} \\
\\ 
\texttt{ANSWER: no} \\
\\ 
\texttt{ <sep> human:} \\
            }
        }
    }
    & & \\ \\

    \theutterance \stepcounter{utterance}  
    & & & \multicolumn{2}{p{0.3\linewidth}}{
        \cellcolor[rgb]{0.9,0.9,0.9}{
            \makecell[{{p{\linewidth}}}]{
                \texttt{\tiny{[GM$|$GM]}}
                \texttt{NA} \\
            }
        }
    }
    & & \\ \\

    \theutterance \stepcounter{utterance}  
    & & & \multicolumn{2}{p{0.3\linewidth}}{
        \cellcolor[rgb]{0.9,0.9,0.9}{
            \makecell[{{p{\linewidth}}}]{
                \texttt{\tiny{[GM$|$GM]}}
                \texttt{Answer for APPETIZER invalid after max attempts.} \\
            }
        }
    }
    & & \\ \\

    \theutterance \stepcounter{utterance}  
    & & & \multicolumn{2}{p{0.3\linewidth}}{
        \cellcolor[rgb]{0.9,0.9,0.9}{
            \makecell[{{p{\linewidth}}}]{
                \texttt{\tiny{[GM$|$GM]}}
                \texttt{Answer is incorrect.} \\
            }
        }
    }
    & & \\ \\

    \theutterance \stepcounter{utterance}  
    & & & \multicolumn{2}{p{0.3\linewidth}}{
        \cellcolor[rgb]{0.9,0.9,0.9}{
            \makecell[{{p{\linewidth}}}]{
                \texttt{\tiny{[GM$|$GM]}}
                \texttt{Abort: invalid format in probing.} \\
            }
        }
    }
    & & \\ \\

\end{supertabular}
}

\end{document}
